% 第 11 题图甲：探究金属丝阻值随气压变化规律的装置（TikZ 线稿）
% 玻璃管内封装金属丝（两端引出接测量电路），抽气泵调节管内气压
\begin{tikzpicture}[font=\small, >=Stealth]
  % 玻璃管（水平，两端密封）
  \draw[thick] (0,1.0) -- (5.4,1.0);
  \draw[thick] (0,1.8) -- (5.4,1.8);
  \draw[very thick] (0,0.95) -- (0,1.85);
  \draw[very thick] (5.4,0.95) -- (5.4,1.85);
  % 金属丝（管内波浪线；弧度写法避免 deg() 维度溢出）
  \draw[very thick, domain=0.35:5.05, samples=100]
    plot (\x, {1.4+0.16*sin(((\x-0.35)*8.021) r)});
  % 金属丝两端引出（接测量电路）
  \draw (0.15,1.4) -- (-1.1,1.4);
  \draw (5.25,1.4) -- (6.2,1.4);
  \node[left] at (-1.1,1.4) {接测量电路};
  \node[below] at (2.7,0.82) {玻璃管};
  \node[above] at (2.7,1.95) {金属丝};
  % 抽气泵（经支管与玻璃管相连）
  \draw[thick] (4.1,1.8) -- (4.1,2.5);
  \draw[thick] (4.1,2.5) -- (5.3,2.5);
  \draw[thick] (5.3,1.9) rectangle (6.7,3.1);
  \node at (6.0,2.5) {抽气泵};
  \draw[->] (6.0,3.1) -- (6.0,3.5) node[above] {抽气};
\end{tikzpicture}
